% ===========================================================================
%  preamble.tex --- packages and macros for the mathematical introduction.
%  When this chapter is dropped into the thesis, merge the contents of this
%  file into the thesis preamble; duplicates may simply be deleted.
% ===========================================================================

\usepackage[T1]{fontenc}
\usepackage[utf8]{inputenc}
\usepackage{lmodern}
\usepackage{microtype}
\usepackage{amsmath,amssymb,amsthm,mathtools}
\usepackage{graphicx}
\usepackage{float}
\usepackage{placeins}
\usepackage{booktabs}
\usepackage[hypcap=false]{caption}
\usepackage{subcaption}
\usepackage{tikz}
\usepackage{pgfplots}
\pgfplotsset{compat=1.18}
\usepackage{enumitem}
\usepackage{xcolor}
\usepackage{array}
\usepackage{hyperref}
\usepackage[nameinlink,noabbrev]{cleveref}

% --- Numbered edit highlights (cut-list review markup) ----------------------
% Full-measure yellow boxes + red number badge.  Toggle off with:
%   \renewcommand{\hln}[2]{#2}\renewcommand{\hlnbox}[2]{#2}
\definecolor{edithl}{RGB}{255,245,150}
\definecolor{editbadge}{RGB}{185,28,28}
\newcommand{\hlnbadge}[1]{%
  \leavevmode
  {\setlength{\fboxsep}{1.1pt}%
   \colorbox{editbadge}{\color{white}\tiny\bfseries\strut\,#1\,}}%
  \nobreak\hspace{0.12em}\relax
}
\newcommand{\hln}[2]{%
  \leavevmode
  \ifhmode\unskip\fi
  \par\nopagebreak
  \noindent
  {\setlength{\fboxsep}{2.5pt}%
   \setlength{\fboxrule}{0pt}%
   \colorbox{edithl}{%
     \parbox{\dimexpr\linewidth-2\fboxsep-1pt\relax}{%
       \raggedright\sloppy
       \hlnbadge{#1}\ignorespaces #2%
     }}}%
  \par
}
\newcommand{\hlnbox}[2]{%
  \par\nopagebreak
  \noindent
  {\setlength{\fboxsep}{4pt}%
   \setlength{\fboxrule}{0pt}%
   \colorbox{edithl}{%
     \parbox{\dimexpr\linewidth-2\fboxsep-1pt\relax}{%
       \raggedright\sloppy
       \hlnbadge{#1}\par\nopagebreak
       #2%
     }}}%
  \par
}

\graphicspath{{figures/}}

% --- Equation cross-references as bare parenthesised numbers -----------------
\crefformat{equation}{(#2#1#3)}
\Crefformat{equation}{(#2#1#3)}
\crefrangeformat{equation}{(#3#1#4)--(#5#2#6)}
\crefmultiformat{equation}{(#2#1#3)}{ and (#2#1#3)}{, (#2#1#3)}{ and (#2#1#3)}

% --- Shared style for the cobweb panels of the Galton--Watson section -------
\pgfplotsset{
    ch2plotstyle/.style={
        width=4.9cm, height=4.9cm,
        axis lines=left,
        xlabel={$S$},
        xmin=0, xmax=1, ymin=0, ymax=1,
        xtick={0,0.5,1}, ytick={0,0.5,1},
        grid=major, grid style={dashed, gray!30},
        samples=100, domain=0:1,
        thick,
        tick label style={font=\tiny},
        label style={font=\small}
    }
}

% --- Notation ---------------------------------------------------------------
% Scalars are set in italic; bold is reserved for vector-valued states, which
% appear only in the multi-type work of later chapters.
\newcommand{\pgf}{probability generating function}
\newcommand{\rv}{random variable}
\newcommand{\pr}[1]{\mathbb{P}\!\left(#1\right)}
\newcommand{\ex}[1]{\mathbb{E}\!\left(#1\right)}
\newcommand{\var}[1]{\operatorname{Var}\!\left(#1\right)}
\newcommand{\dd}{\mathrm{d}}
\newcommand{\ee}{\mathrm{e}}
\newcommand{\ddt}[1]{\frac{\dd #1}{\dd t}}
\newcommand{\dda}[2]{\frac{\dd #1}{\dd #2}}
\newcommand{\pddt}[1]{\frac{\partial #1}{\partial t}}
\newcommand{\pdd}[2]{\frac{\partial #1}{\partial #2}}
\newcommand{\lrb}[1]{\left(#1\right)}
\newcommand{\lrbs}[1]{\left[#1\right]}
\newcommand{\Ac}{A_{\mathrm{c}}}
\newcommand{\ind}{\mathbf 1}

% --- Cross-chapter references -----------------------------------------------
% Redefining these two macros is all that thesis integration requires: they
% become "Chapter~\ref{ch:...}" once the chapter order is fixed.
\newcommand{\ChM}{the mathematical introduction}
\newcommand{\ChA}{this chapter}
\newcommand{\ChAcap}{This chapter}

% --- Forward references to later modelling chapters -------------------------
% Every forward reference is routed through \fwd, whose first argument is a
% stable key and whose second is the prose actually typeset.
\newcommand{\fwd}[2]{#2}

% --- Theorem environments ---------------------------------------------------
\theoremstyle{plain}
\newtheorem{theorem}{Theorem}[chapter]
\newtheorem{proposition}[theorem]{Proposition}
\newtheorem{lemma}[theorem]{Lemma}
\newtheorem{corollary}[theorem]{Corollary}
\theoremstyle{definition}
\newtheorem{definition}[theorem]{Definition}
\theoremstyle{remark}
\newtheorem{remark}[theorem]{Remark}

\allowdisplaybreaks
